\section{Reflexion \& Fazit}

\subsection{Was lief gut}

Ein wesentlicher Erfolgsfaktor war die schrittweise Entwicklung über klar
abgrenzte Meilensteine.
Dadurch konnten Anforderungen iterativ umgesetzt und laufend validiert werden:
zuerst ein stabiles UI-Grundgerüst, danach die Migration auf eine
komponentenbasierte SPA und schließlich die Erweiterung zur Full-Stack-Lösung
mit Backend, Authentifizierung und Datenbank.
Diese Reihenfolge war rückblickend sinnvoll, weil sie technische Risiken
reduzierte und zu jedem Zeitpunkt einen lauffähigen Projektstand sicherstellte.

Auch die Verbindung von fachlicher Domäne und technischer Umsetzung hat gut
funktioniert.
Da das Projekt aus einem realen Bedarf im eigenen Spielkontext entstanden ist,
waren Anforderungen wie Übersichtlichkeit, schnelle Bedienbarkeit und
rollenbasierte Sichtbarkeit von Beginn an klar.
Das hat Priorisierungsentscheidungen erleichtert und zu einer konsistenten
Produktlinie geführt.

In der Teamarbeit war positiv, dass technische Umbauten nicht nur
feature-getrieben, sondern auch strukturell durchgeführt wurden.
Beispiele dafür sind die Umstellung auf wiederverwendbare Komponenten, die
Einführung von TypeScript-Typen, die spätere Trennung in Frontend/Backend und
die klare Dokumentation der Meilensteine.
So wurde vermieden, dass kurzfristige Lösungen die Wartbarkeit dauerhaft
verschlechtern.

Zusätzlich war die frühe Etablierung von Tooling hilfreich:
Linting/Formatierung, automatisierte Tests und ein reproduzierbarer
Startprozess mit Docker Compose haben die Qualität in der Abschlussphase
messbar stabilisiert.
Gerade bei der Vorbereitung von Demo und Ausarbeitung hat sich gezeigt, dass
ein konsistenter Betriebsstand viele Folgekosten reduziert.

\subsection{Was würden wir anders machen}

Rückblickend würde das Projekt einige technische Entscheidungen früher treffen.
Die wichtigste davon betrifft die Systemgrenze zwischen Frontend und Backend.
In den frühen Ständen wurden Teile der Logik zunächst clientseitig modelliert;
mit wachsender Komplexität (Rollen, Authentifizierung, persistente Inhalte)
mussten diese Konzepte später auf serverseitige Regeln umgestellt werden.
Beim nächsten Projekt würde die API- und Datenmodellplanung früher erfolgen, um
Umbauten in späteren Phasen zu reduzieren.

Ein zweiter Punkt betrifft die Benennung und Konsistenz von Pfaden, Modulen und
Routen.
Durch den langen Entwicklungsverlauf mit mehreren Umstrukturierungen
(\texttt{src/}~$\rightarrow$~\texttt{frontend/}, neue Backend-Struktur,
Refactorings im Routing) entstand zwischenzeitlich inkonsistente Terminologie.
Diese wurde zwar bereinigt, hätte aber mit früheren Konventionen weniger
Nacharbeit erzeugt.

Auch die Teststrategie würde künftig früher breiter aufgesetzt.
Die vorhandenen Tests decken zentrale Pfade gut ab, wurden aber erst in späteren
Meilensteinen systematisch ausgebaut.
Ein früheres, stärker risikobasiertes Testdesign (z.\,B. Auth- und
Rollenfälle bereits parallel zur ersten Backend-Integration) hätte manche
Fehlerkorrekturen beschleunigt.

Organisatorisch würde die Abschlussphase zeitlich noch stärker gegen Ende
gepuffert werden.
Der inhaltliche Abschluss (Demo, Kapitelabgleich, formale Nacharbeit) ist
erfahrungsgemäß aufwendiger als erwartet, selbst wenn die Kernfunktionalität
bereits steht.
Für zukünftige Projekte ist deshalb eine frühere Synchronisation zwischen
Implementierungsstand und finalem Abgabeformat sinnvoll.

\subsection{Was wurde gelernt}

Technisch wurde vor allem deutlich, wie wichtig eine saubere
Verantwortungstrennung in Webanwendungen ist.
Die Trennung zwischen UI, API und Persistenz verbessert nicht nur die
Wartbarkeit, sondern auch Sicherheit und Nachvollziehbarkeit.
Insbesondere bei Authentifizierung, Rollenlogik und Datenzugriff ist eine
serverseitige Durchsetzung zentral, da rein clientseitige Regeln nicht
ausreichen.

Ein weiterer Lernpunkt ist der praktische Wert konsistenter Typisierung.
TypeScript hat die Kommunikation zwischen Modulen verbessert und Refactorings
abgesichert, vor allem bei komplexeren Datenstrukturen für Kampagneninhalte und
Benutzerzustände.
In Kombination mit klaren Interfaces und einem strukturierten API-Vertrag
konnten Integrationsprobleme früh erkannt werden.

Im Bereich Qualitätssicherung wurde gelernt, dass Tests den größten Nutzen
haben, wenn sie reale Kernabläufe abbilden:
Login/Logout, geschützte Routen, Rollenrechte und typische Nutzeraktionen.
Diese Tests sind nicht nur technische Absicherung, sondern auch eine
ausführbare Spezifikation des erwarteten Systemverhaltens.

Organisatorisch hat das Projekt gezeigt, dass kontinuierliche Dokumentation ein
entscheidender Erfolgsfaktor ist.
Wenn Architektur-, Betriebs- und Testentscheidungen erst am Ende beschrieben
werden, geht Kontext verloren.
Die laufende Pflege von README, Meilensteinzuordnung und Ausarbeitung erleichtert
dagegen sowohl Teamabstimmung als auch Abgabevorbereitung erheblich.

Insgesamt wurde aus dem Projekt mitgenommen, dass ein gutes Ergebnis weniger von
einzelnen Features abhängt als von der Verbindung aus
\textbf{inkrementeller Planung, strukturiertem Refactoring, testbarer
Architektur und reproduzierbarem Betrieb}.
Diese Kombination hat aus einem anfänglich kleinen Frontend-Prototyp eine
vollständig nutzbare Full-Stack-Anwendung gemacht.
